\chapter{Calcium}
\index{Calcium}

\section*{Definition and Overview}
Calcium is a mineral that is essential for health. More than 99\% of the body's calcium is stored in bones and teeth, where it provides strength and structure. The small amount present in blood and body fluids is vital for muscle contraction, nerve signalling, blood clotting, and the normal function of the heart.

Because these functions are essential, the body tightly controls blood calcium levels, drawing on bone stores when dietary intake is low. Calcium is obtained from food and, where needed, from supplements. Disorders of calcium are usually problems of hormonal regulation or bone health rather than of diet alone, and both deficiency and excess can cause significant illness.

\section*{Epidemiology}
Calcium is needed throughout life, and intake varies widely between countries and dietary patterns. Many adults, particularly older women, do not obtain enough calcium from food, and low intakes are common in people who avoid dairy products. Osteoporosis, which is linked to long-term calcium deficiency, affects roughly one in ten people over the age of 50 and is considerably more common in women. Despite low dietary intakes, most people maintain normal blood calcium levels because the body draws on bone reserves, which is why blood levels are a poor measure of overall calcium status.

\section*{Aetiology and Pathophysiology}
The body keeps blood calcium within a narrow range using hormones -- parathyroid hormone (PTH) and vitamin D -- that adjust absorption from the gut, reabsorption by the kidneys, and the release or storage of calcium in bone:
\begin{itemize}
    \item \textbf{Dietary intake:} dairy foods, leafy green vegetables, nuts, seeds, and fortified foods provide calcium; long-term low intake weakens bones
    \item \textbf{Vitamin D:} vitamin D is required for calcium absorption, so deficiency can cause low calcium despite adequate intake
    \item \textbf{Hormone control:} parathyroid hormone and vitamin D raise blood calcium, while calcitonin lowers it
    \item \textbf{Bone turnover:} bone constantly releases and takes up calcium, acting as a reserve that protects blood levels
    \item \textbf{Imbalance:} too much or too little calcium in the blood usually reflects a problem with hormones, kidneys, or vitamin D rather than diet alone
\end{itemize}

\section*{Clinical Features}
Calcium itself causes no symptoms in health, but abnormal blood levels produce distinctive features:
\begin{itemize}
    \item \textbf{Low blood calcium (hypocalcaemia):} muscle cramps and twitching, tingling in the fingers and around the mouth, fatigue, and, in severe cases, seizures
    \item \textbf{High blood calcium (hypercalcaemia):} thirst, frequent urination, constipation, nausea, weakness, confusion, and kidney stones
    \item \textbf{Low dietary calcium over time:} few symptoms until bone loss becomes advanced
    \item \textbf{Deficiency in children:} slowed growth, poor tooth development, and, rarely, rickets
    \item \textbf{Bone disease:} fragile bones, back pain, loss of height, and fractures as late consequences of prolonged calcium shortage
\end{itemize}

\section*{Differential Diagnosis}
Symptoms and findings related to calcium can reflect other conditions:
\begin{itemize}
    \item Muscle cramps and tingling from magnesium deficiency, thyroid problems, or nerve conditions
    \item Fatigue and weakness from anaemia, underactive thyroid, or chronic illness
    \item Bone pain and fractures from injury, vitamin D deficiency, or bone diseases other than osteoporosis
    \item Thirst and frequent urination from diabetes
    \item Kidney stones from causes other than high calcium, such as gout or dehydration
\end{itemize}

\section*{When to Seek Medical Care}
A blood calcium level outside the normal range is usually detected by a doctor, and any abnormal result needs review.

\textbf{Contact your local emergency services for:}
\begin{itemize}
    \item Severe muscle cramping, twitching, or a seizure
    \item Confusion, drowsiness, or loss of consciousness
    \item Chest pain or an irregular heartbeat associated with an abnormal calcium level
\end{itemize}

Otherwise, people should see a doctor if they have persistent thirst, bone pain, muscle cramps, or other unexplained symptoms, or if a recent blood test showed abnormal calcium.

\section*{Diagnosis and Investigations}
Calcium problems are diagnosed through blood tests and assessment of related minerals and hormones:
\begin{itemize}
    \item \textbf{Blood test:} total and sometimes ionised (free) calcium in the blood
    \item \textbf{Albumin:} measured together with calcium, because albumin binds calcium and affects the total result
    \item \textbf{Vitamin D:} checked because deficiency is a common cause of low calcium
    \item \textbf{Parathyroid hormone (PTH):} helps determine why calcium is too high or too low
    \item \textbf{Magnesium and kidney function:} measured because both affect calcium regulation
    \item \textbf{Bone density scan (DEXA):} used to assess bone strength when osteoporosis is suspected
    \item \textbf{Urine tests:} may be used to check how much calcium the body is losing
\end{itemize}

\section*{Management}
Management depends on whether the problem lies in diet, bone health, or abnormal blood levels:
\begin{itemize}
    \item \textbf{Eat calcium-rich foods:} dairy products, sardines and tinned fish with bones, leafy greens, tofu, and fortified plant milks
    \item \textbf{Ensure adequate vitamin D:} sunlight exposure and, when advised, vitamin D supplements support calcium absorption
    \item \textbf{Supplements:} calcium tablets may be recommended when dietary intake is inadequate, especially for older adults, at modest doses
    \item \textbf{Treat the cause:} correcting a low or high blood calcium usually means treating the underlying hormonal, kidney, or vitamin D problem
    \item \textbf{Medicines for bone health:} for diagnosed osteoporosis, specific medicines may be prescribed in addition to calcium and vitamin D
    \item \textbf{Follow-up:} repeat blood tests confirm that calcium has returned to a normal level
\end{itemize}

\section*{Complications}
\begin{itemize}
    \item Osteoporosis and fractures after long-term calcium shortage
    \item Kidney stones, which can be triggered by high calcium intake in some people
    \item Constipation and stomach upset with large calcium supplements
    \item Interactions, as calcium supplements can reduce the absorption of some medicines, such as iron and thyroid medication
    \item Serious complications of severe abnormal blood levels, including seizures, abnormal heart rhythms, or coma
\end{itemize}

\section*{Prognosis and Outcomes}
With adequate calcium and vitamin D, bone health can usually be maintained, and low dietary intakes are easy to correct. Abnormal blood calcium levels often resolve once the underlying cause is treated, and severe cases generally improve with hospital care. For osteoporosis, treatment slows further bone loss and reduces fracture risk, although calcium alone does not repair bone that is already lost. In most people, a balanced diet provides all the calcium the body needs.

\section*{Prevention and Self-Care}
\begin{itemize}
    \item Eat calcium-rich foods daily, with the amount depending on age, sex, and life stage
    \item Get regular, safe sunlight exposure and discuss vitamin D with your doctor if you are at risk of deficiency
    \item Do weight-bearing exercise, which helps keep bones strong
    \item Avoid smoking and limit alcohol, as both weaken bone
    \item Do not take calcium supplements without advice, since excess can be harmful
    \item Take calcium supplements at a different time from iron or thyroid medicines
    \item Have a bone density test if you are in a high-risk group, such as post-menopausal women or people on long-term steroid treatment
\end{itemize}

\section*{Sources and Further Reading}
The following sources provide reliable, accessible information on calcium and bone health:
\begin{itemize}
    \item NHS. \textit{Vitamins and minerals: calcium.} https://www.nhs.uk/conditions/vitamins-and-minerals/
    \item National Institutes of Health, Office of Dietary Supplements. \textit{Calcium fact sheet.} https://ods.od.nih.gov/
    \item Mayo Clinic. \textit{Calcium and calcium supplements.} https://www.mayoclinic.org/
    \item MedlinePlus. \textit{Calcium.} https://medlineplus.gov/calcium.html
    \item Osteoporosis Australia. \textit{Calcium and bone health.} https://www.osteoporosis.org.au/
    \item MSD Manual. \textit{Calcium disorders.} https://www.msdmanuals.com/
\end{itemize}
